% Vietnamese translation for OI Public Library
% Source-File: IOI中国国家候选队论文/2003/周源/周源.ppt
\documentclass[11pt]{article}
\usepackage{oipl-vn}

\title{Bản trình chiếu: Biểu diễn nhỏ nhất của chuỗi vòng}
\author{Zhou Yuan}
\date{2003}

\begin{document}
\maketitle

\origtitle{Slide companion for minimal representation of cyclic strings}
\sourcepath{IOI China candidate team papers / 2003 / Zhou Yuan / Zhou Yuan.ppt}

\section{Phạm vi bản dịch}

Bản trình chiếu này đi kèm bài viết về biểu diễn nhỏ nhất trong bài toán đẳng
cấu vòng của chuỗi. Các slide khôi phục được chứa định nghĩa phép quay
\(s^{(k)}\), ví dụ \texttt{babba}, \texttt{bbaba}, và các so sánh độ phức
tạp \(O(N^2)\), \(O(N)\).

\section{Định nghĩa}

Với chuỗi \(s\), ký hiệu \(|s|\) là độ dài. Phép quay một bước đưa ký tự đầu
về cuối:
\[
  s^{(1)} = s[2..|s|] + s[1],
\]
và \(s^{(k)}\) là phép quay lặp lại \(k\) lần. Hai chuỗi có cùng độ dài là
đẳng cấu vòng nếu một chuỗi là một phép quay của chuỗi kia.

\section{Biểu diễn nhỏ nhất}

Cách đơn giản là xét toàn bộ \(s,s^{(1)},\ldots,s^{(|s|-1)}\), mất
\(O(N^2)\). Cách tốt hơn là nối \(S=s+s\) và tìm vị trí bắt đầu của biểu diễn
từ điển nhỏ nhất trong thời gian \(O(N)\). Khi đó hai chuỗi vòng tương đương
khi và chỉ khi biểu diễn nhỏ nhất của chúng giống nhau.

\section{Ứng dụng}

Tư tưởng chọn đại diện chuẩn cũng áp dụng cho trạng thái tìm kiếm và kiểm tra
đẳng cấu: các đối tượng tương đương được gom về cùng một biểu diễn để tránh
xử lý lặp.

\end{document}
